Ze loerde om haar heen
en draaide op één been

tot ze naar struiken keek.
Haar gezicht werd lijkbleek.

Emmy zag heidampen.
Had ze soms te kampen

met een stille voyeur?
Open ging de hutdeur.

Wat ruisde er kriskras
vlug door het struikgewas?

Ze versnelde haar pas
naar een nabij moeras.

De jacht was begonnen.
De inkomstenbronnen

kozen als het mis zat
vaker het hazepad.

De jager kende het
en had een val gezet

ergens in het rietgras
rondom het veenmoeras.

De strik strak gespannen
wist te overmannen

menig zoogdier als aas.
Nu was Emmy de haas.

Ze werd zonder geweld
door de jager geveld.

Ze was totaal verward.
Haar gevoelige hart

was op hol geslagen.
Haar borstkas moest dragen

de stokkende adem
net onder haar boezem.

Haar mond was droog en naar
en haar tong voelde zwaar.

Haar handen waren klam
en haar knieën zwak lam.

Een knoop lag in haar maag.
Haar spieren leken traag.

De angst greep zich vast.
Haar gevoel werd verkast

en bleef onderbewust.
Een paniek dat niet sust

daar het niet werd erkend
of zelfs niet werd herkend.

“Zou ze worden geslacht?”
was waar Emmy aan dacht?

Het mes in zijn schede
kon met één halssnede

al haar bloeddoorlopen
vlot laten weglopen.

Het nare paasmenu
passeerde de revue.

Weer klonk het jammeren
van al die lammeren.

Nog maar 6 maanden oud
maakte men hen al koud.

Of het haar overkwam,
ook al was zij geen lam,

proefde Emmy de vrees
opgeslorpt door hun vlees.

Ze had meer met dieren.
Die konden niet klieren.

Er waaide vooral stof
toen een jonge Schwarzkopf,

direct na het zogen,
vlak voor Emmy’s ogen
gevoelloos werd gekeeld.
Een onuitwisbaar beeld.

Het ondierlijke kwaad
telde in het kwadraat

op met scheiding misplaatst
en in zichzelf weerkaatst.

De schuine zichtzijde
liet de assen snijden

in vegetarisme.
De wortel van rede

gaf de exacte maat
van haar goedige staat.

Omdat ze moreel was
net als Pythagoras

kon ze geen dier eten.
Plantaardig diëten

was toen vrij bijzonder
en daar viel ook onder

Emmy haar omgeving.
Dus werd haar weigering

om een lijk van een dier
te nuttigen, alhier

waar men vee consumeert,
geenszins gerespecteerd.

Emmy mijdde stiekem
elk maal van het schorem

dat ze moest verslinden
als ze iets kon vinden

wat een dier was voordien.
Johannes vers 14.
Het woord werd immers vlees
was Emmy grootste vrees.

Geen scheiding tussen God
en wat in de stoofpot

op zondag was gekookt
of apart was gerookt.

Deze moraliteit
paste de vrome meid

ook toe op alcohol.
Mensen werden er dol

en zondig van, zag ze.
Geen alcoholisme

was echt een uitdaging
omdat je het ontving

en men het zelfs opdrong,
ook al was je nog jong,

als bron van energie.
Echter niet voor Emmy.

Dit dieet hield haar slank.
Haar reehuid bleef roomblank.

Haar huid was handelswaar.
De jager was dankbaar.

Haar enige lompen
lagen ingekrompen

ver buiten haar bereik.
Volledige inkijk

gaf ze haar kidnapper.
Hij meende haar knapper

dan elke rosse vrouw,
die hij in ogenschouw

tot zich had genomen.
Hij ging door met bomen:

“Wie heeft dit lief wezen
hardvochtig verwezen

naar dit ruige bestaan?
Wat heeft ze nou misdaan?

Veel edele heren
zullen haar begeren.

Haar willen bevrijden
en zich zo verblijden

door met haar te trouwen.
Wat heeft haar weerhouden?”

Vele malen groter
en ook zoveel sterker

dan Emmy bracht hij haar
polsen dicht bij elkaar

vlak boven haar hoofdje.
Het opkijkend roofje

oogde zo machteloos,
terwijl ze zich roerloos

gaf aan de inspectie.
Hij oversteeg Emmy

met drie decimeter.
Anderhalve meter

bedroeg haar lengtemaat.
Blikken uit zijn gelaat

verkennen haar lichaam.
Ze probeert behoedzaam

zijn zicht te onttrekken
door dat te bedekken

waarvoor gold in die tijd
als een intimiteit.

Kort na een worsteling
tegen de gijzeling

stevig in zijn polsgreep,
was het dat ze begreep

dat verzet zinloos was.
Ze gaf geen tegengas

toen hij van top tot teen
zijn licht over haar scheen.

Emmy was nat en koud.
Het bleek het aanmaakhout

voor adrenaline
alsook endorfine.

Er ontstond opwinding
uit deze ontbering.

“Kom dichter bij Christus”,
zo zei Sint-Franciscus:

“en draag je eigen kruis”.
Emmy voelde zich thuis

in pijn of ongemak,
waarin loutering stak.

Haar huid speelde van streek
verstoppertje zo bleek,

want hij werd blootgesteld.
Alles werd braaf verteld.

Terwijl hij haar lichaam
zorgvuldig en langzaam

deel voor deel verkende,
kwam het zo bekende

gevoel van “zijn geveld”
als ze werd blootgesteld.

Tenminste in Emmy
haar rijke fantasie.

Ze stelde zich vroeger,
als knecht van een ploeger,

hem als sergeant-majoor
wel eens verbeeldend voor.

Als een stoere soldaat
met een dwingend gelaat.

Toen ze fantaseerde
dat hij haar chanteerde

om al haar werkkleding,
voor een overtreding

van zijn eigen krijgswet,
te geven als debet,

dan ervoer ze prikkels.
Er waren meer bikkels

die haar zo opwonden.
Die haar boeien konden.

De vreemde opwinding
die ze weer onderging.

Door deze sterke drang
ging de jager zijn gang

zonder enig verzet.
Emmy verlangde het.

Haar voedsel-omhulsel
werd, zonder haar vulsel

met het Goddelijke,
het verrukkelijke.

Emmy haar snijwond
was waar hij het licht vond

naar Emmy’s openheid
en Rumi’s kwetsbaarheid.

Vloog zijn kwieke eros
naar haar microkosmos?
